Bahasa Isyarat Indonesia (BISINDO) merupakan bahasa alami yang digunakan oleh komunitas Tuli di Indonesia
dan memiliki karakteristik gerakan, ekspresi, serta konteks yang kaya.
Dalam praktiknya, terdapat sejumlah kosakata BISINDO yang secara visual tampak serupa,
namun memiliki perbedaan makna yang ditentukan oleh arah gerakan, tempo,
posisi, serta unsur non-manual.

Seiring berkembangnya teknologi, muncul berbagai upaya pemanfaatan sistem berbasis komputer
untuk membantu pemahaman bahasa isyarat.
Namun, sebagian pendekatan yang ada masih berfokus pada aspek teknis semata
dan belum sepenuhnya melibatkan komunitas Tuli sebagai mitra keilmuan.
Penelitian ini berangkat dari kebutuhan untuk mengkaji pola gerakan BISINDO
secara etis, kolaboratif, dan kontekstual.

Penelitian ini tidak bertujuan menggantikan peran guru, penerjemah,
atau praktik pembelajaran BISINDO yang telah ada.
Fokus penelitian diarahkan pada kajian pola gerakan isyarat
sebagai bagian dari pengembangan sistem bantu berbasis teknologi
dalam konteks akademik,
dengan pendampingan dan validasi dari PUSBISINDO Jawa Barat.
